\documentclass[11pt, a4paper]{article}

% ── Packages ──────────────────────────────────────────────────────────
\usepackage[utf8]{inputenc}
\usepackage[T1]{fontenc}
\usepackage{lmodern}
\usepackage[margin=1in]{geometry}
\usepackage{amsmath, amssymb, amsthm}
\usepackage{booktabs}
\usepackage{longtable}
\usepackage{graphicx}
\usepackage{hyperref}
\usepackage{xcolor}
\usepackage{listings}
\usepackage{tcolorbox}
\usepackage{polya}

\hypersetup{
  colorlinks=true,
  linkcolor=polyablue,
  urlcolor=polyablue,
  citecolor=polyablue,
}

% ── Code listing style ────────────────────────────────────────────────
\lstdefinestyle{polya-python}{
  language=Python,
  basicstyle=\ttfamily\small,
  keywordstyle=\color{polyablue}\bfseries,
  commentstyle=\color{polyagray}\itshape,
  stringstyle=\color{polyagreen},
  numbers=left,
  numberstyle=\tiny\color{polyagray},
  numbersep=8pt,
  frame=single,
  framerule=0.4pt,
  rulecolor=\color{polyagray},
  breaklines=true,
  showstringspaces=false,
  tabsize=4,
}
\lstset{style=polya-python}
\title{Land Survey Analysis}
\author{uber-polya}
\date{2026-02-26}

\begin{document}

\maketitle
\tableofcontents
\newpage

% ── Phase A: Problem Formulation ─────────────────────────────────────
\section{Problem Formulation}

\begin{polyamodel}

\textbf{Problem Type}: Find \hfill
\textbf{Domain}: Computational Geometry


\end{polyamodel}

% ── Universe ──────────────────────────────────────────────────────────
\subsection{Universe}

\begin{itemize}
  \item $Building Footprint: \{width, height, center_x, center_y\}$
  \item $Max Inscribed Rectangle: \{center_x, center_y, width, height, area_m2\}$
\end{itemize}

% ── Variables ─────────────────────────────────────────────────────────
\subsection{Variables}

\begin{longtable}{lll}
\toprule
\textbf{Symbol} & \textbf{Meaning} & \textbf{Type / Range} \\
\midrule
\endhead
$x$ & decision variable & $\mathbb{{R}}$ \\
\bottomrule
\end{longtable}

% ── Structure ─────────────────────────────────────────────────────────
\subsection{Structure}

A surveyor has recorded GPS coordinates for 8 boundary markers defining an irregular lot. Given these coordinates (in meters, local projected coordinate system), compute:

1. **Area** of the lot polygon
2. **Perimeter** of the lot boundary
3. **Convex hull** of the survey points
4. **Centroid** (geometric center of mass)
5. **Building fit check** -- determine whether a proposed rectangular building footprint fits entirely within the lot

The lot is an irregular 8-sided polygon. The building footprint is an axis-aligned rectangle placed at a proposed location. The solver must verify containment using independent geometric methods.

% ── Mapping ───────────────────────────────────────────────────────────
\subsection{Real-World Mapping}

\begin{longtable}{ll}
\toprule
\textbf{Real-World Concept} & \textbf{Mathematical Object} \\
\midrule
\endhead
problem input & $instance parameters$ \\
\bottomrule
\end{longtable}

% ── Constraints ───────────────────────────────────────────────────────
\subsection{Constraints}

\begin{enumerate}
  \item $Computational geometry$
  \hfill \textit{(algorithms operating on points, lines, and polygons in the plane)}
  \item $Polygon operations$
  \hfill \textit{(area, perimeter, convex hull, centroid for arbitrary simple polygons)}
  \item $Spatial containment$
  \hfill \textit{(testing whether one geometric object lies entirely within another)}
  \item $Voronoi partitioning$
  \hfill \textit{(dual of Delaunay triangulation, useful for nearest-facility analysis)}
  \item $Shoelace formula$
  \hfill \textit{(exact O(n) area computation for simple polygons from vertex coordinates)}
  \item $Convex hull$
  \hfill \textit{(smallest convex set containing all points; measures lot "regularity")}
\end{enumerate}

% ── Objective / Claim ─────────────────────────────────────────────────
\subsection{Claim}


% ── Approach ──────────────────────────────────────────────────────────
\subsection{Approach}

\begin{description}
  \item[Suggested approach:] Shoelace + ConvexHull(Qhull) + Shapely containment
  \item[Complexity class:] O(n)
  \item[Available tools:] Shoelace + ConvexHull
\end{description}
% ── Phase B: Solution ────────────────────────────────────────────────
\section{Solution}

\begin{polyaresult}

\textbf{Answer}: Solution computed


\textbf{Optimal}: No / Unknown \hfill
\textbf{Feasible}: Yes
\textbf{Algorithm}: Shoelace + ConvexHull(Qhull) + Shapely containment \quad $O(n)$

\textbf{Time}: 1.864808s


\end{polyaresult}

% ── Solution Details ──────────────────────────────────────────────────
\subsection{Solution Details}


Method: 1. Area: Shoelace formula (Gauss's area formula) for simple polygons -- O(n)
2. Perimeter: Sum of Euclidean distances between consecutive vertices -- O(n)
3. Convex hull: Scipy ConvexHull (Qhull library, O(n log n))
4. Centroid: Weighted average using the shoelace sub-expressions -- O(n)
5. Containment: Shapely Polygon.contains() for the full building rectangle -- O(n) per query
6. Max inscribed rectangle: Grid search over candidate positions and sizes within the polygon

% ── Verification ──────────────────────────────────────────────────────
\subsection{Verification}

\begin{longtable}{lcl}
\toprule
\textbf{Check} & \textbf{Status} & \textbf{Detail} \\
\midrule
\endhead
Feasibility & \passmark & All constraints satisfied \\
\bottomrule
\end{longtable}

% ── Phase C: Interpretation ──────────────────────────────────────────
\section{Interpretation}

% ── The Question & The Answer ─────────────────────────────────────────
\subsection{The Question}
A surveyor has recorded GPS coordinates for 8 boundary markers defining an irregular lot.

\subsection{The Answer}

\begin{polyainsight}[title={Bottom Line}]
A feasible solution was found.
\end{polyainsight}

% ── What This Means ───────────────────────────────────────────────────
\subsection{What This Means}

- Lot area computed via shapely and independently verified with the shoelace formula
- Perimeter of the lot boundary
- Convex hull vertices and hull area (compared to lot area for concavity measure)
- Centroid coordinates
- Building containment verdict (fits / does not fit)
- Maximum axis-aligned inscribed rectangle estimate
- Independent verification of all computed quantities

1. Area: Shoelace formula (Gauss's area formula) for simple polygons -- O(n)
2. Perimeter: Sum of Euclidean distances between consecutive vertices -- O(n)
3. Convex hull: Scipy ConvexHull (Qhull library, O(n log n))
4. Centroid: Weighted average using the shoelace sub-expressions -- O(n)
5. Containment: Shapely Polygon.contains() for the full building rectangle -- O(n) per query
6. Max inscribed rectangle: Grid search over candidate positions and sizes within the polygon

% ── Sensitivity Analysis ──────────────────────────────────────────────

% ── Figures ───────────────────────────────────────────────────────────

% ── Recommendations ───────────────────────────────────────────────────
\subsection{Recommendations}

\begin{enumerate}
  \item Review the model assumptions to ensure real-world fidelity.
\end{enumerate}

% ── Limitations ───────────────────────────────────────────────────────
\subsection{Limitations}

\begin{polyawarnbox}
\begin{itemize}
  \item Model assumes deterministic parameters; real-world variability not captured.
  \item Solution quality depends on the accuracy of input data.
\end{itemize}
\end{polyawarnbox}

% ── Appendix ─────────────────────────────────────────────────────────

\end{document}